\section{Lexical Structure}

\subsection*{Whitespace and Comments}
Whitespace separates tokens. Line comments begin with \texttt{;} and continue to
end of line. Block comments are delimited by \texttt{\#|} and \texttt{|\#} and
may span multiple lines. Datum comments use \texttt{\#;} and discard the next
syntactic datum.

\subsection*{Delimiters}
Parentheses \texttt{( )} and square brackets \texttt{[ ]} are interchangeable and
can be freely mixed, as long as they are properly balanced.

\subsection*{Identifiers}
Identifiers are any non-whitespace sequences that are not one of the structural
punctuation characters \texttt{( ) [ ] ; |}. This means identifiers may contain
characters such as \texttt{+ - * / < > = ? !} and may begin with digits. Numeric
lexemes are recognized separately (see Section \emph{Numbers}).

\subsection*{Booleans}
Booleans are written as \texttt{\#t} and \texttt{\#f}.

\subsection*{Strings}
Strings are delimited by double quotes \texttt{"..."} and may contain any
character except an unescaped double quote. Newlines are permitted inside
strings and are preserved.

\subsection*{Quoting Tokens}
Shorthand quoting tokens are supported: \texttt{'} (quote), \texttt{`} (quasiquote),
\texttt{,} (unquote), and \texttt{,@} (unquote-splicing).

\subsection*{Vectors}
Vector literals use the reader form \texttt{\#( ... )} or \texttt{\#[ ... ]}.
